\documentclass[10pt,a4paper]{amsart}
\usepackage{amsmath,amssymb,amsthm,graphicx}
% Retained minimal preamble for the Stacks Project chapter more-algebra.tex (More on Algebra).
% The chapter file pinned-original-source.tex carries only \input{preamble} on its first line, so no
% preamble text lives inside the chapter itself.  The Stacks Project uses one single shared preamble for
% every chapter; the snapshot stages it identically as sources/stacks-chow/preamble.tex and
% sources/stacks-etale-cohomology/preamble.tex.  The block below is copied mechanically, line for line,
% from sources/stacks-chow/preamble.tex lines 14-18 (the Xy-pic setup) and 154-217 (the native theorem
% declarations together with the mathematical macros).  Nothing was added, removed or redefined; the
% wrapper class and the amsmath/amssymb/amsthm/graphicx base packages are supplied by the assembler.
% Chapter: more-algebra.tex -- More on Algebra.
% For commutative diagrams we use Xy-pic
\usepackage[all]{xy}

% We use 2cell for 2-commutative diagrams.
\xyoption{2cell}

% Theorem environments.
%
\theoremstyle{plain}
\newtheorem{theorem}[subsection]{Theorem}
\newtheorem{proposition}[subsection]{Proposition}
\newtheorem{lemma}[subsection]{Lemma}

\theoremstyle{definition}
\newtheorem{definition}[subsection]{Definition}
\newtheorem{example}[subsection]{Example}
\newtheorem{exercise}[subsection]{Exercise}
\newtheorem{situation}[subsection]{Situation}

\theoremstyle{remark}
\newtheorem{remark}[subsection]{Remark}
\newtheorem{remarks}[subsection]{Remarks}

\numberwithin{equation}{subsection}

% Macros
%
\def\lim{\mathop{\mathrm{lim}}\nolimits}
\def\colim{\mathop{\mathrm{colim}}\nolimits}
\def\Spec{\mathop{\mathrm{Spec}}}
\def\Hom{\mathop{\mathrm{Hom}}\nolimits}
\def\Ext{\mathop{\mathrm{Ext}}\nolimits}
\def\SheafHom{\mathop{\mathcal{H}\!\mathit{om}}\nolimits}
\def\SheafExt{\mathop{\mathcal{E}\!\mathit{xt}}\nolimits}
\def\Sch{\mathit{Sch}}
\def\Mor{\mathop{\mathrm{Mor}}\nolimits}
\def\Ob{\mathop{\mathrm{Ob}}\nolimits}
\def\Sh{\mathop{\mathit{Sh}}\nolimits}
\def\NL{\mathop{N\!L}\nolimits}
\def\CH{\mathop{\mathrm{CH}}\nolimits}
\def\proetale{{pro\text{-}\acute{e}tale}}
\def\etale{{\acute{e}tale}}
\def\QCoh{\mathit{QCoh}}
\def\Ker{\mathop{\mathrm{Ker}}}
\def\Im{\mathop{\mathrm{Im}}}
\def\Coker{\mathop{\mathrm{Coker}}}
\def\Coim{\mathop{\mathrm{Coim}}}

% Boxtimes
%
\DeclareMathSymbol{\boxtimes}{\mathbin}{AMSa}{"02}

%
% Macros for moduli stacks/spaces
%
\def\QCohstack{\mathcal{QC}\!\mathit{oh}}
\def\Cohstack{\mathcal{C}\!\mathit{oh}}
\def\Spacesstack{\mathcal{S}\!\mathit{paces}}
\def\Quotfunctor{\mathrm{Quot}}
\def\Hilbfunctor{\mathrm{Hilb}}
\def\Curvesstack{\mathcal{C}\!\mathit{urves}}
\def\Polarizedstack{\mathcal{P}\!\mathit{olarized}}
\def\Complexesstack{\mathcal{C}\!\mathit{omplexes}}
% \Pic is the operator that assigns to X its picard group, usage \Pic(X)
% \Picardstack_{X/B} denotes the Picard stack of X over B
% \Picardfunctor_{X/B} denotes the Picard functor of X over B
\def\Pic{\mathop{\mathrm{Pic}}\nolimits}
\def\Picardstack{\mathcal{P}\!\mathit{ic}}
\def\Picardfunctor{\mathrm{Pic}}
\def\Deformationcategory{\mathcal{D}\!\mathit{ef}}
\usepackage{cleveref}
\crefname{originallemma}{original lemma}{original lemmas}
\Crefname{originallemma}{Original lemma}{Original lemmas}
\makeatletter
\newlabel{lemma-element-projective}{{0GVH}{source}{}{sourcepointer-1}{}}
\newlabel{lemma-element-projective@cref}{{[originallemma][0][]0GVH}{[1][0][]source}{}{}{}}
\makeatother
\crefname{originallemma}{original lemma}{original lemmas}
\Crefname{originallemma}{Original lemma}{Original lemmas}
\makeatletter
\newlabel{lemma-trick-to-find-good-element}{{0GVI}{source}{}{sourcepointer-2}{}}
\newlabel{lemma-trick-to-find-good-element@cref}{{[originallemma][0][]0GVI}{[1][0][]source}{}{}{}}
\makeatother
\title{Every element of a projectively infinite-rank module over a ring with Noetherian radical quotient lies in a finite stably free direct summand}
\author{The Stacks Project Authors}
\date{stacks-project:more-algebra.tex:local-snapshot-20261012; GNU FDL 1.2 or later}
\begin{document}
\maketitle
\setcounter{subsection}{1}
\noindent\textit{Editorial note.} All delivered mathematical text is the author's own, copied from the pinned Stacks Project chapter file more-algebra.tex; no statement or proof step was rewritten, reordered or removed. The counted statement is the labeled lemma environment at source lines 39107-39114 (canonical tag 0GVJ) and the counted proof is the proof environment at source lines 39116-39150, which immediately follows it and is its own proof.\par\medskip
\noindent\textit{Editorial note.} The proof contains one footnote at source lines 39134-39141 verifying that the projection onto K'' is surjective. It is the author's own text inside the counted span and is retained verbatim.\par\medskip
\noindent\textit{Editorial note.} The two lemmas the proof invokes, lemma-element-projective and lemma-trick-to-find-good-element, are earlier statements of the same section of this chapter and are carried as navigation stubs with their roles recorded in the reasons.\par\medskip
\noindent\textit{Editorial note.} The chapter file pinned-original-source.tex carries only an input of the shared preamble on its first line, so no preamble text lives inside the chapter. The retained file sources/stacks-more-algebra/retained-preamble.tex is a mechanical, line-for-line copy of the project's single shared preamble as staged in the snapshot (Xy-pic setup plus the native theorem declarations and mathematical macros). Nothing was added, removed or redefined, and the excerpt serves only as preamble, never as a span coordinate.\par\medskip
\noindent\textit{Editorial note.} The source URL is the source-level Stacks Project URL, as required when more than one card is produced from one source; the per-claim tag is carried by claim\_key and by this note. The counted statement and proof cite no bibliography entry, so this card needs no bibliography.\par\medskip
\noindent\textit{Allowed standard prerequisites (editorial).}\begin{itemize}
\item An element of a projective module that is not a zero divisor on the fibres generates a free direct summand after one passes to a suitable element of the module. invoked at proof lines 39126-39129, where the surjection onto the complement K' is used to find p with the image of (1, p) generating a free direct summand. Produces the free summand that is split off from the complement inside R direct sum P.
\item A projective module that is the kernel of a surjection onto a projective module is a direct summand, and being a direct summand of a projective module, hence projective. used at proof lines 39142-39144. Makes the kernel of the projection onto K'' a direct summand of P containing the given element s.
\item Stable freeness is preserved by taking direct summands of finite free modules: if R direct sum N is isomorphic to a finite free module then N is stably free. used at proof lines 39144-39149. Concludes that the direct summand produced is finite stably free, which is the conclusion of the lemma.
\end{itemize}\medskip
\ifdefined\hypertarget\hypertarget{sourcepointer-1}{}\fi
\noindent\textit{Source reference.} Original lemma 0GVH: sources/stacks-more-algebra/pinned-original-source.tex, source lines 39009--39009. The reference supplies the auxiliary statement that places (0, s) inside a finite free direct summand of a finite free direct summand of R direct sum P, used at proof lines 39117-39120; the reduction and everything after it are inside the counted span.\par\medskip
\ifdefined\hypertarget\hypertarget{sourcepointer-2}{}\fi
\noindent\textit{Source reference.} Original lemma 0GVI: sources/stacks-more-algebra/pinned-original-source.tex, source lines 39026--39026. The reference supplies the element p whose image generates a free direct summand of the complement, used at proof lines 39127-39130; the decomposition of the complement and the stable freeness conclusion are inside the counted span.\par\medskip
\begin{lemma}
\label{lemma-element-in-free-summand}
Let $R$ be a ring with Jacobson radical $J$ such that $R/J$
is Noetherian. Let $P$ be a projective $R$-module
such that $P_\mathfrak m$ has infinite rank for all maximal ideals
$\mathfrak m$ of $R$. Let $s \in P$. Then we can find
a finite stably free direct summand $M \subset P$ such that $s \in M$.
\end{lemma}

\begin{proof}
By Lemma \ref{lemma-element-projective} we can find a finite free
module $F$ and a finite free direct summand $K \subset F \oplus P$ such
that $(0, s) \in K$. By induction on the rank of $F$ we reduce
to the case discussed in the next paragraph.

\medskip\noindent
Assume there exists a finite stably free direct summand $K \subset R \oplus P$
such that $(0, s) \in K$. Choose a complement $K'$ of $K$, i.e.,
such that $R \oplus P = K \oplus K'$. The projection
$\pi : R \oplus P \to K'$ is surjectve, hence by
Lemma \ref{lemma-trick-to-find-good-element}
we find a $p \in P$ such that $\pi(1, p) \in K'$ generates
a free direct summand. Accordingly we write
$K' = R\pi(1, p) \oplus K''$. We see that
$$
R \oplus P = K \oplus K' = K \oplus R\pi(1, p) \oplus K''
$$
The projection $\pi' : P \to K''$ is surjective\footnote{Namely,
if $k'' \in K''$ then $k''$ viewed as an element of $K'$
can be written as $k'' = \lambda \pi(1, 0) + \pi(0, q)$
for some $\lambda \in R$ and $q \in P$.
This means $k'' = \lambda \pi(1, p) + \pi(0, q - \lambda p)$.
This in turn means that $q - \lambda p$ maps to $k''$ by
the composition $P \to R \oplus P \xrightarrow{\pi} K' \to K''$
since $K' \to K''$ annihilates $\pi(1, p)$.}
and hence split
(as $K''$ is projective). Thus $\Ker(\pi') \subset P$ is a
direct summand containing $s$. Finally, by construction we
have an isomorphism
$$
R \oplus \Ker(\pi') \cong K \oplus R\pi(1, p)
$$
and hence since $K$ is finite and stably free, so is $\Ker(\pi')$.
\end{proof}

\end{document}
